\scanpage{9}\mspage{4}
fixes qui porte la structure\ldots). On peut aussi décrire ces structures en termes de $S$ ou de $A$ par la donnée d'un groupe $R$ (``forme'' de $\mathbb{Z}/4\mathbb{Z}$) opérant sur $S$ resp. sur $A$, qui devient un torseur sous $\Gamma$
\[
 S\text{ torseur sous }\Gamma\tag*{(2 bis)}
\]
\[
 A\text{ torseur sous }\Gamma\tag*{(3 bis)}
\]
ce qui fait finalement \deleted{quatre} \inserted{cinq} descriptions de \deleted{\unsure{la}} $=$ structure de ``carré''. Je veux expliciter les passages de l'une de ces descriptions à l'autre.

$\left.\begin{aligned}2\text{ bis})&\longmapsto2)\\3\text{ bis})&\longmapsto3)\end{aligned}\right\}$
en prenant l'unique élément $\underline a$ d'ordre $2$ de $\Gamma$, qui est une \unsure{d'involution|d'une involution|involution} ss pts fixes de $S$ (resp $A$), et on pose
\[
D=S/\underline a,\quad\text{resp. }C=A/\underline a
\]

$2)$\deleted{bis}$\implies1)$ \quad Soit $A\subset\mathcal{P}_2(S)\inserted{\subset\mathcal{P}(S)}$ défini par $A=\mathcal{P}_2(S)\setminus\mathrm{Im}\,D$, c'est un ens. à $4$ éléments, \deleted{qui} et on a une relation tautologique $R\subset S\times A$

$3)\implies1)$ \quad Soit $S\subset\mathcal{P}_2(A)$ défini par $S=\mathcal{P}_{\unsure{2}}(A)\setminus\mathrm{Im}\,C$, ens à $4$ éléments et on a relation tautologique $R'\subset A\times S$ d'où $ {}^{s}R'\stackrel{\mathrm{déf}}{=}R\subset S\times A$

$\left.\begin{aligned}1&\implies2\text{ bis}\\1&\implies3\text{ bis}\end{aligned}\right)$
On considère $\unsure{\mathfrak{D}|\mathscr{D}}=\mathrm{Aut}\underbrace{Q}_{Q=\{S,A,R\}}$
\diagram{\begin{tikzcd}[ampersand replacement=\&,row sep=small]
\&\unsure{\mathfrak{S}|G}_S\\
\deleted{\unsure{\mathfrak{D}_Q|\mathfrak{S}_Q|\mathfrak{D}}}\arrow[ur,hook]\arrow[dr,hook]\&\\
\&\mathfrak{S}_A
\end{tikzcd}}
\deleted{\unsure{on a|ou}} \deleted{\unsure{de} $\mathfrak{D}\unsure{\cap|{}_{A},} V_Q$} \inserted{ordre $8$} \inserted{2-sous-gr. de Sylow de $\mathfrak{S}_S$ ou $\mathfrak{S}_A$}
L'image inverse de $V_S\subset\mathfrak{S}_S$ et celle de $V_A\subset\mathfrak{S}_A$ \deleted{\unsure{sont}} deux $=$ sous-groupes d'ordre $4$ \deleted{\unsure{distincts|d'ordre}} de $\mathfrak{D}$, notés $V_{\unsure{S|Q}}$ et $V_A$, isomorphes respectivement à $\mathbb{F}_2^C$ et $\mathbb{F}_2^D$ -- leur intersection est leur sous

\deleted{(\unsure{de $S$, $A$}) \unsure{dans} $\unsure{\mathfrak{D}}$ \illegible{2} en un groupe $V_Q$ qui \illegible{1} vectoriel de rg $2$ sur $\mathbb{F}_2$, il y a deux éléments d'ordre $4$ \unsure{distincts} inverses l'un de l'autre, engendrant}
\diagram{\begin{tikzcd}[ampersand replacement=\&,column sep=small,row sep=small]
\mathbb{F}_2^C\simeq V_S\arrow[d,hook]\arrow[r,hook]\&\mathfrak{S}_S\\
\mathfrak{D}\arrow[ur,hook,no head]\arrow[dr,hook]\&\\
\mathbb{F}_2^D\simeq V_A\arrow[u,hook]\arrow[r,hook]\&\mathfrak{S}_A
\end{tikzcd}}
% archive sheet number: 8
